\documentclass[a4paper]{ltjsarticle}
% \usepackage{luatexja}
\usepackage{amsmath}
\usepackage{amssymb}
\usepackage{amsfonts}
\usepackage{braket}
\usepackage{bm}
\usepackage{geometry}
% \usepackage{newcomputermodern}
% Hiraginoフォントがシステムにインストールされている場合、以下のコメントを解除してください
% \usepackage{luatexja-fontspec}
% \setmainjfont[UprightFont=・{Hiragino Mincho ProN W3}, BoldFont={Hiragino Mincho ProN W6}]{Hiragino Mincho ProN}
% \setsansjfont[UprightFont={Hiragino Kaku Gothic ProN W3}, BoldFont={Hiragino Kaku Gothic ProN W6}]{Hiragino Kaku Gothic ProN}

\title{Peskin2章}
\author{}
\date{\today}

\begin{document}

\maketitle

- p.14 U(t) の計算
時間並進の演算子は$e^{-i H t}$なので、$x_0$から$x$に遷移する確率振幅$U(t)$は
\begin{align*}
  U(t) & = \bra{x} e^{-i H t} \ket{x_0} \\
       & = \bra{x} e^{-i p^2 / (2m)} \ket{x_0} \\
       & = \int \frac{d^3p}{(2\pi)^3} \bra{x} e^{-i p^2 / (2m)} \ket{p} \braket{p|x_0} & \text{（完全性関係を利用）} \\
       & = \int \frac{d^3p}{(2\pi)^3} e^{-i p^2 / (2m)} \braket{x|p} \braket{p|x_0} \\
       & = \int \frac{d^3p}{(2\pi)^3} e^{-i p^2 / (2m)} e^{i \bm{p} \cdot (\bm{x} - \bm{x}_0)} & (\braket{x|p} = e^{i \bm{x} \cdot \bm{p}}) \\
       & = \int \frac{d^3p}{(2\pi)^3} e^{- \frac{it}{2m}(p + \frac{m}{t}(\bm{x} - \bm{x}_0))^2} e^{\frac{im}{2t}(\bm{x} - \bm{x}_0)^2} \\
       & = \frac{1}{(2\pi)^3} \left( \frac{2m\pi}{it} \right)^{3/2} e^{\frac{im}{2t}(\bm{x} - \bm{x}_0)^2} & \text{（ガウス積分）} \\
       & = \left( \frac{m}{2\pi i t} \right)^{3/2} e^{\frac{im}{2t}(\bm{x} - \bm{x}_0)^2}
\end{align*}
4行目では運動量の固有値を取るので、指数関数はブラケットの外に出すことができる。

- p.14二つ目の式の計算
\begin{align*}
   & \int d^3p \, e^{-it\sqrt{p^2 + m^2}} e^{i \bm{p} \cdot (\bm{x} - \bm{x}_0)} \\
   & = \int_0^{2\pi} d\phi \int_0^\infty p^2 dp \int_{-1}^1 d\mu \, e^{-it\sqrt{p^2 + m^2}} e^{ip|\bm{x} - \bm{x}_0|\mu} \\
   & = 2\pi \int_0^\infty p^2 dp \frac{1}{ip|\bm{x} - \bm{x}_0|} e^{-it\sqrt{p^2 + m^2}} (e^{ip|\bm{x} - \bm{x}_0|} - e^{-ip|\bm{x} - \bm{x}_0|}) \\
   & = 4\pi \int_0^\infty dp \, p \sin(p|\bm{x} - \bm{x}_0|) e^{-it\sqrt{p^2 + m^2}}
\end{align*}

- (2.2), (2.3)式の導出
\[
  \partial_\mu \left(\frac{\partial \mathcal{L}}{\partial(\partial_\mu \phi)} \delta \phi\right) = \partial_\mu \left(\frac{\partial \mathcal{L}}{\partial(\partial_\mu \phi)}\right) \delta \phi + \frac{\partial \mathcal{L}}{\partial(\partial_\mu \phi)} \partial_\mu (\delta \phi), \quad \partial_\mu (\delta \phi) = \delta(\partial_\mu \phi)
\]
であるから、
\begin{align*}
  0 & = \delta S \\
    & = \int d^4 x \, \delta \mathcal{L} \\
    & = \int d^4 x \left\{ \frac{\partial \mathcal{L}}{\partial \phi}\delta \phi + \frac{\partial \mathcal{L}}{\partial(\partial_\mu \phi)}\delta(\partial_\mu \phi) \right\} \\
    & = \int d^4 x \left\{ \frac{\partial \mathcal{L}}{\partial \phi}\delta \phi - \partial_\mu \left(\frac{\partial \mathcal{L}}{\partial(\partial_\mu \phi)}\right) \delta \phi + \partial_\mu \left(\frac{\partial \mathcal{L}}{\partial(\partial_\mu \phi)} \delta \phi\right) \right\} & \text{（上の計算を代入）}
\end{align*}
作用積分が収束するには場$\phi$が遠方では急速に0に近づく必要があるから、無限遠の値を参照する表面積分の項は0となる。また、変分$\delta \phi$は任意であるから、変分法の原理より$\delta \phi$の係数は常に0である。したがって
\[
  \frac{\partial \mathcal{L}}{\partial \phi} - \partial_\mu \left(\frac{\partial \mathcal{L}}{\partial(\partial_\mu \phi)}\right) = 0
\]

- Noerther's Theorem の計算

作用積分Sが
\[
  S = \int d^4 x \, \mathcal{L}
\]
で与えられているので、何らかの変分に対しラグランジアン密度が
\[
  \mathcal{L} \to \mathcal{L}' = \mathcal{L} + \alpha \partial_\mu \mathcal{J}^\mu (x)
\]
（$\alpha$は微小定数）と変化する限りは作用積分は変わらず、したがって運動方程式も不変である。
\begin{align*}
  \because S' & = \int d^4 x \, \mathcal{L}' \\
              & = \int d^4 x (\mathcal{L} + \alpha \partial_\mu \mathcal{J}^\mu (x)) \\
              & = S + \alpha \int d^4 x \, \partial_\mu \mathcal{J}^\mu (x)
\end{align*}
となるが表面積分は先に見たように0となるからである。

ここで変分
\[
  \phi(x) \to \phi'(x) = \phi(x) + \alpha \Delta \phi(x)
\]
において、$\mathcal{L}$の変化を見てみよう。先の計算と同様にして
\[
  \alpha \Delta \mathcal{L} = \frac{\partial \mathcal{L}}{\partial \phi} (\alpha \Delta \phi) + \left(\frac{\partial \mathcal{L}}{\partial(\partial_\mu \phi)}\right)\partial_\mu (\alpha \Delta \phi)
  = \alpha \partial_\mu \left(\frac{\partial \mathcal{L}}{\partial(\partial_\mu \phi)} \Delta \phi\right) + \alpha \left[\frac{\partial \mathcal{L}}{\partial \phi} - \partial_\mu \left(\frac{\partial \mathcal{L}}{\partial(\partial_\mu \phi)}\right)\right] \Delta \phi
\]
と書けるが、第二項はオイラー方程式から0となる。したがって
\[
  j^\mu (x) := \frac{\partial \mathcal{L}}{\partial(\partial_\mu \phi)} \Delta \phi - \mathcal{J}^\mu
\]
と定めると
\[
  \partial_\mu j^\mu (x) = 0
\]
となる。

- (2.15), (2.16)式の導出
\[
  \phi \to e^{i\alpha} \phi = (1 + i\alpha + O(\alpha^2)) \phi
\]
なので$\alpha$が微少量である時、2次の項を無視して(2.15)式を得る。

また
\[
\mathcal{L} = |\partial_\mu \phi|^2 - m^2 |\phi|^2
\]
のとき、この変分でラグランジアンは変化しないから、
\[
  \mathcal{J} = 0
\]
であり、
\[
  j^\mu = \frac{\partial \mathcal{L}}{\partial(\partial_\mu \phi)} \Delta \phi + \text{c.c.}
\]
を計算することで(2.15)式を得る。（ここでc.c.はComplex conjugate, 複素共役の略で、直前の項の複素共役である。今回$\phi$は複素場なので$\phi^*$による微分の項が出てくることに注意）

- (2.24)式の導出
\[
  \phi = \frac{1}{\sqrt{2\omega}} (a + a^\dagger), \quad p = -i \sqrt{\frac{\omega}{2}}(a - a^\dagger)
\]
と
\[
  [\phi, p] = i
\]
より
\begin{align*}
  i & = [\phi, p] \\
    & = - \frac{i}{2}  [(a + a^\dagger), (a - a^\dagger)] \\
    & = -\frac{i}{2} \{[a^\dagger, a] - [a, a^\dagger]\} \\
    & = i [a, a^\dagger]
\end{align*}
となる。したがって
\[
[a, a^\dagger] = 1
\]

- (2.25), (2.26)式の意味
クラインゴルドン場を単振動場と同様にフーリエ空間で表したいが、エネルギーが波数に依存することから分かるように、生成消滅演算子$a, a^\dagger$も波数に依存する。そこで、$\phi, \pi$が実場であることを考慮すると、
\begin{align*}
  \phi(x) & = \int \frac{d^3\bm{p}}{(2\pi)^3} \frac{1}{\sqrt{2\omega_{\bm{p}}}} ( a_{\bm{p}} e^{i \bm{p} \cdot \bm{x}} + a_{\bm{p}}^\dagger e^{-i \bm{p} \cdot \bm{x}}) \\
  \pi(x)  & = \int \frac{d^3\bm{p}}{(2\pi)^3} (-i) \sqrt{\frac{\omega_{\bm{p}}}{2}} ( a_{\bm{p}} e^{i \bm{p} \cdot \bm{x}} - a_{\bm{p}}^\dagger e^{-i \bm{p} \cdot \bm{x}})
\end{align*}
と表せるのである。ただし、意味ありげに出てきた$\frac{1}{\sqrt{2\omega_{\bm{p}}}}$と$(-i)\sqrt{\frac{\omega_{\bm{p}}}{2}}$はこれらが相対論的に不変である（ローレンツ変換で変化しない）ために必要。

- (2.30)の計算
\begin{align*}
  \phi(x) & = \int \frac{d^3\bm{p}}{(2\pi)^3} \frac{1}{\sqrt{2\omega_{\bm{p}}}} ( a_{\bm{p}} + a_{-\bm{p}}^\dagger ) e^{i \bm{p} \cdot \bm{x}} \\
  \pi(x)  & = \int \frac{d^3\bm{p}'}{(2\pi)^3} (-i) \sqrt{\frac{\omega_{\bm{p}'}}{2}} ( a_{\bm{p}'} - a_{-\bm{p}'}^\dagger )e^{i \bm{p}' \cdot \bm{x}}
\end{align*}
と$\pi(x)$の積分変数を$p'$と変えてから代入し、交換関係
\[
  [a_{\bm{p}}, a^\dagger_{\bm{p}'}] = (2\pi)^3 \delta^{(3)}(\bm{p} - \bm{p}'), \quad [a_{\bm{p}}, a_{\bm{p}'}] = [a^\dagger_{\bm{p}}, a^\dagger_{\bm{p}'}] = 0
\]
を使って整理すると
\begin{align*}
  [\phi(x), \pi(x')]
   & = \int \frac{d^3\bm{p} \, d^3\bm{p}'}{(2\pi)^6} \frac{-i}{2} \sqrt{\frac{\omega_{\bm{p}'}}{\omega_{\bm{p}}}} ([a^\dagger_{-\bm{p}}, a_{\bm{p}'}] + [a_{\bm{p}}, -a^\dagger_{-\bm{p}'}]) e^{i(\bm{p} \cdot \bm{x} + \bm{p}' \cdot \bm{x}')} \\
   & = \int \frac{d^3\bm{p} \, d^3\bm{p}'}{(2\pi)^6} \frac{-i}{2} \sqrt{\frac{\omega_{\bm{p}'}}{\omega_{\bm{p}}}} \{ (-1)(2\pi)^3 \delta^{(3)}(\bm{p} + \bm{p}') - (2\pi)^3 \delta^{(3)}(\bm{p} + \bm{p}') \} e^{i(\bm{p} \cdot \bm{x} + \bm{p}' \cdot \bm{x}')} \\
   & = \int \frac{d^3\bm{p}}{(2\pi)^3} i \sqrt{\frac{\omega_{-\bm{p}}}{\omega_{\bm{p}}}} e^{i \bm{p}(\bm{x} - \bm{x}')} \\
   & = \frac{i}{(2\pi)^3} \int d^3\bm{p} \, e^{i \bm{p}(\bm{x} - \bm{x}')} \\
   & = i \delta^{(3)}(\bm{x} - \bm{x}')
\end{align*}
ここで、$\omega_{\bm{p}} = \omega_{-\bm{p}}$であることに注意

- (2.31)の計算
これも(2.30)の計算と同様に$\phi, \pi$の片方の積分変数を$p'$に変えてから代入すればよく、
\[
  \nabla \phi = \int \frac{d^3\bm{p}}{(2\pi)^3} \frac{i\bm{p}}{\sqrt{2\omega_{\bm{p}}}} ( a_{\bm{p}} + a_{-\bm{p}}^\dagger ) e^{i \bm{p} \cdot \bm{x}}
\]
であるから
\begin{align*}
  H & = \int d^3x \left[ \frac{1}{2}\pi^2 + \frac{1}{2}(\nabla \phi)^2 + \frac{1}{2}m^2 \phi^2 \right] \\
    & = \int d^3x \int \frac{d^3\bm{p} \, d^3\bm{p}'}{(2\pi)^6} \\
    & \quad e^{i(\bm{p}+\bm{p}') \cdot \bm{x}} \left\{
  \frac{1}{2} \frac{\sqrt{\omega_{\bm{p}} \omega_{\bm{p}'}}}{2} (a_{\bm{p}} - a^\dagger_{-\bm{p}})(a_{\bm{p}'} - a^\dagger_{-\bm{p}'}) + \frac{1}{2} \frac{-\bm{p} \cdot \bm{p}' + m^2}{2\sqrt{\omega_{\bm{p}} \omega_{\bm{p}'}}} (a_{\bm{p}} + a^\dagger_{-\bm{p}})(a_{\bm{p}'} + a^\dagger_{-\bm{p}'})
  \right\}
\end{align*}
ここで
\[
  \int d^3x \, e^{i(\bm{p}+\bm{p}') \cdot \bm{x}} = (2\pi)^3 \delta(\bm{p}+\bm{p}')
\]
であるから, $x$について積分した後$p'$について積分し、$\omega_{\bm{p}} = \omega_{-\bm{p}}$と$p^2 + m^2 = \omega_p^2$であることに注意すると
\begin{align*}
  H
   & = \int \frac{d^3\bm{p}}{(2\pi)^3} \frac{\omega_{\bm{p}}}{4} [-(a_{\bm{p}} - a_{-\bm{p}}^\dagger)(a_{-\bm{p}} - a_{\bm{p}}^\dagger) + (a_{\bm{p}} + a_{-\bm{p}}^\dagger)(a_{-\bm{p}} + a_{\bm{p}}^\dagger)] \\
   & = \int \frac{d^3\bm{p}}{(2\pi)^3} \frac{\omega_{\bm{p}}}{4} \cdot 2(a_{\bm{p}} a_{\bm{p}}^\dagger + a_{-\bm{p}}^\dagger a_{-\bm{p}})
\end{align*}
第二項目の積分変数を$\bm{p} \to -\bm{p}$と変えれば、
\begin{align*}
  H & = \int \frac{d^3\bm{p}}{(2\pi)^3} \frac{\omega_{\bm{p}}}{2} (a_{\bm{p}} a_{\bm{p}}^\dagger + a_{\bm{p}}^\dagger a_{\bm{p}}) \\
    & = \int \frac{d^3\bm{p}}{(2\pi)^3} \omega_{\bm{p}} (a_{\bm{p}} a_{\bm{p}}^\dagger + \frac{1}{2}[a_{\bm{p}}, a_{\bm{p}}^\dagger])
\end{align*}

- (2.31)の計算
\[
  [AB, C] = A[B, C] + [A, C]B
\]
を使う。

(2.30)式の第二項のデルタ関数を無視すると
\[
  [H, a_{\bm{p}}^\dagger] = \int \frac{d^3\bm{p}'}{(2\pi)^3} \omega_{\bm{p}'} [a_{\bm{p}'} a_{\bm{p}'}^\dagger, a_{\bm{p}}^\dagger]
\]
\[
  [a_{\bm{p}'} a_{\bm{p}'}^\dagger, a_{\bm{p}}^\dagger] = a_{\bm{p}'} [a_{\bm{p}'}^\dagger, a_{\bm{p}}^\dagger] + [a_{\bm{p}'}, a_{\bm{p}}^\dagger] a_{\bm{p}'}^\dagger = (2\pi)^3 \delta(\bm{p}' - \bm{p}) a_{\bm{p}'}^\dagger
\]
なので
\[
  [H, a_{\bm{p}}^\dagger] = \omega_{\bm{p}} a_{\bm{p}}^\dagger
\]

- (2.33)の計算
\begin{align*}
  \bm{P} & = - \int d^3x \, \pi \nabla \phi \\
             & = -\int \frac{d^3\bm{p} \, d^3\bm{p}'}{(2\pi)^3} \delta(\bm{p}+\bm{p}') \frac{\bm{p}}{2} \sqrt{\frac{\omega_{\bm{p}}}{\omega_{\bm{p}'}}} (a_{\bm{p}} - a_{-\bm{p}}^\dagger)(a_{\bm{p}'} + a_{-\bm{p}'}^\dagger) \\
             & = - \int \frac{d^3\bm{p}}{(2\pi)^3} \frac{\bm{p}}{2} (a_{\bm{p}} a_{-\bm{p}} + a_{\bm{p}} a_{\bm{p}}^\dagger - a_{-\bm{p}}^\dagger a_{-\bm{p}} - a_{-\bm{p}}^\dagger a_{\bm{p}}^\dagger)
\end{align*}
ここで、$p a_p a_{-p}$と$p a_{-p}^\dagger a_p^\dagger$は奇関数であるから、積分すると消える。したがって
\[
  \bm{P} = - \int \frac{d^3\bm{p}}{(2\pi)^3} \frac{\bm{p}}{2} (a_{\bm{p}} a_{\bm{p}}^\dagger - a_{-\bm{p}}^\dagger a_{-\bm{p}})
\]
また第二項は$\bm{p} \to -\bm{p}$と積分変数を変更すると
\begin{align*}
  \bm{P} & = \int \frac{d^3\bm{p}}{(2\pi)^3} \frac{\bm{p}}{2} (a_{\bm{p}}^\dagger a_{\bm{p}} - a_{\bm{p}} a_{\bm{p}}^\dagger) \\
             & = \int \frac{d^3\bm{p}}{(2\pi)^3} \bm{p} (a_{\bm{p}}^\dagger a_{\bm{p}} + \frac{1}{2}[a_{\bm{p}}, a_{\bm{p}}^\dagger])
\end{align*}
最後の交換関係から出てくるデルタ関数はp.22上部の議論から無視できるので、(2.33)式を得る。

\end{document}
